\documentclass[10pt]{article}
\usepackage[margin=0.6in]{geometry}
\usepackage{amsmath,amssymb}
\usepackage{multicol}
\usepackage{booktabs}
\usepackage{array}
\usepackage{enumitem}
\usepackage[dvipsnames]{xcolor}

\pagestyle{empty}

% Tableau signs
\newcommand{\T}{\mathbb{T}}
\newcommand{\F}{\mathbb{F}}

% Vertical dots shorthand
\newcommand{\vmark}{\,\vert\,}

% Section formatting
\newcommand{\ruleheader}[1]{\medskip\noindent\textbf{\large #1}\par\smallskip\hrule\smallskip}
\newcommand{\rulesubheader}[1]{\smallskip\noindent\textbf{#1}\par\smallskip}

% Branching notation
\newcommand{\branches}[2]{%
  \begin{tabular}{@{}c@{\qquad}c@{}}
  #1 & #2
  \end{tabular}%
}

\setlength{\parindent}{0pt}
\setlength{\parskip}{2pt}

\begin{document}

\begin{center}
{\LARGE\bfseries Practice Midterm}\\[4pt]
{\large Phil 305 --- Propositional and Modal Logic}
\end{center}

\ruleheader{Part 1: S5}

Use the simplified rules for \textbf{S5} to build tableaux for each of these claims. Say which of them are valid, i.e., are theorems of \textbf{S5}.

\begin{enumerate}
\item $\Box(\Box p \rightarrow (q \vee r)) \vee \Box(\Box q \rightarrow p)$
\item $\Diamond \Box (p \rightarrow q) \rightarrow (\Box p \rightarrow \Diamond q)$
\item $\Box (p \vee q) \rightarrow (\Diamond (\neg q \vee r) \rightarrow \Diamond (p \vee r))$
\item $(\neg \Diamond p \vee \neg \Diamond \Diamond p) \rightarrow \neg \Diamond \Diamond \Diamond p$
\end{enumerate}

\ruleheader{Part 2: K}

Use the rules for \textbf{K} to build tableaux for each of these claims. Say which of them are valid, i.e., are theorems of \textbf{K}.

\begin{enumerate}
\setcounter{enumi}{4}
\item $(\Diamond p \wedge \Box q) \rightarrow \Diamond q$
\item $\Box (p \wedge \neg p) \vee \Diamond (q \vee \neg q)$
\item $\Box \neg(p \vee \neg q) \rightarrow \Box p$
\end{enumerate}

\ruleheader{Part 3: S4}

Use the rules for \textbf{S4} to build tableaux for each of these claims. Say which of them are valid, i.e., are theorems of \textbf{S4}. Remember \textbf{S4} is the logic \textbf{KT4}, so it has the rules for \textbf{T} and \textbf{4} as well as \textbf{K}

\begin{enumerate}
\setcounter{enumi}{7}
\item $\Box(p \rightarrow (q \wedge r)) \rightarrow (\Diamond \Diamond p \rightarrow \Diamond (q \vee s))$
\item $(\Diamond \Box p \wedge \Diamond \Box q) \rightarrow \Diamond(p \wedge q)$
\item $ \Box \Diamond \Box \Diamond p \rightarrow \Box \Diamond p$
\end{enumerate}

\end{document}